\bookchapter{บทที่ 35:เด็กกับโทรศัพท์—คู่มือการศึกษาในยุคดิจิทัล}

\chapterintro{บทนี้จะไม่กล่าวว่า “สมาร์ทโฟนเป็นภัยพิบัติ” และไม่ส่งเสริมว่า “ยิ่งสัมผัสเร็วยิ่งฉลาด” เราจะแบ่งตามช่วงอายุในการจัดอุปกรณ์ ตั้งกฎเกณฑ์ สอนสื่อสาร และปกป้องสุขภาพ ทำทีละขั้นตอน}

\booksection{ให้มือถือกับลูกตอนอายุเท่าไหร่?}

ไม่มีคำตอบตายตัว แต่มีเกณฑ์อายุเป็นฐาน ถ้าให้เร็วเกินไป เปลือกสมองส่วนหน้าของเด็กยังพัฒนาไม่เต็มที่ ความสามารถในการควบคุมตัวเองจะสู้การออกแบบเย้ายวนของอัลกอริทึมไม่ได้ แต่ถ้าให้ช้าเกินไป อาจทำให้เขาขาดการเชื่อมต่อทางสังคมกับเพื่อนฝูง ด้านล่างนี้คือคำแนะนำแบ่งตามช่วงอายุ:

\begin{itemize}
  \item \textbf{3-6 ปี}:พยายามอย่าให้เด็กแตะมือถือ หากจำเป็นต้องใช้ ควรใช้ครั้งละไม่เกิน 15 นาที เนื้อหาควรเป็นแอนิเมชันเพื่อการศึกษาที่ผู้ปกครองคัดสรรแล้วเป็นหลัก และต้อง\textbf{รับชมด้วยกัน}ในช่วงวัยนี้หน้าจอส่งผลกระทบต่อการมองเห็นและการพัฒนาสมาธิมากที่สุด
  \item \textbf{7-9 ปี}:สามารถให้ยืมใช้งานบนอุปกรณ์ของผู้ปกครองได้ เพื่อใช้เรียนออนไลน์ ฟังนิทาน หรือค้นหาข้อมูล เริ่มสร้างกฎระเบียบ "ต้องขออนุญาตก่อนใช้ และคืนเมื่อใช้เสร็จ" โดยยังไม่ต้องจัดมือถือเฉพาะตัวให้
  \item \textbf{10-12 ปี}:อาจพิจารณาให้โทรศัพท์แบบพื้นฐานหรือนาฬิกาสำหรับเด็ก เพื่อใช้ติดต่อสื่อสารตอนไป-กลับโรงเรียน หากจำเป็นต้องใช้สมาร์ทโฟนจริงๆ ควรเปิดใช้งานระบบควบคุมโดยผู้ปกครองแบบเข้มงวด
  \item \textbf{13 ปีขึ้นไป}:เข้าสู่ระดับมัธยมต้น ความต้องการทางสังคมเพิ่มมากขึ้น สามารถให้สมาร์ทโฟนได้ แต่ต้องใช้ควบคู่กับ "ข้อตกลงการใช้งาน" และการทบทวนเป็นระยะ
\end{itemize}

\begin{tcbinfo}{หลักการ "ใช้พอเพียง"}
มือถือเครื่องแรกของลูกไม่จำเป็นต้องเป็นรุ่นท็อป ประสิทธิภาพที่เหลือเฟือของหน่วยประมวลผลมักถูกใช้ไปกับการเล่นเกม การเลือกมือถือราคาพันกว่าบาทหรือเครื่องตกแต่งใหม่จากโรงงานที่แบตเตอรี่อึด ประสิทธิภาพปานกลาง และหน้าจอถนอมสายตา กลับจะช่วยให้ควบคุมได้ง่ายกว่า งบประมาณอ้างอิง: ช่วง 1500-2500 หยวน ก็เพียงพอสำหรับความต้องการของนักเรียนส่วนใหญ่

\end{tcbinfo}

\booksubsection{นาฬิกาอัจฉริยะเป็นทางเลือกในช่วงเปลี่ยนผ่าน}

หากแค่ต้องการแก้ปัญหาว่า “ติดต่อลูกได้ไหมหลังเลิกเรียน” นาฬิกาสำหรับเด็กจะเหมาะสมกว่ามือถือ เพราะมันโทรได้ ส่งข้อความเสียงได้ และระบุตำแหน่งได้เท่านั้น ไม่มีเบราว์เซอร์และคลิปวิดีโอสั้น ซึ่งเป็นการจำกัด “ช่องทางการเสพติด” ได้ตามธรรมชาติ:

\begin{itemize}
  \item \textbf{นาฬิกาโทรศัพท์ Xiaoai}:ระบบนิเวศครบครันที่สุด การเชื่อมต่อสังคมกับเพื่อนวัยเดียวกันด้วยการ "แตะเพื่อเพิ่มเพื่อน" เป็นจุดขายหลัก แต่ระบบนิเวศ App เฉพาะนั้นค่อนข้างปิด
  \item \textbf{นาฬิกาเด็ก Huawei}:เชื่อมต่อกับมือถือ Huawei ได้ดี ระบุตำแหน่งแม่นยำ รองรับการโทรแบบวิดีโอ คุ้มค่าสูง
  \item \textbf{นาฬิกาเด็ก Xiaomi}:ราคาเป็นมิตร ฟังก์ชันพื้นฐานครบครัน เหมาะสำหรับครอบครัวที่มีงบประมาณจำกัด
\end{itemize}

ข้อควรระวังเมื่อซื้อนาฬิกา: เลือกรุ่นที่รองรับ\textbf{4G ทุกเครือข่าย}(เครือข่าย 2G/3G กำลังปิดให้บริการทีละน้อย) และต้องตรวจสอบอัตราค่าบริการแพ็กเกจสำหรับเด็กของผู้ให้บริการเครือข่าย

\booksection{คู่มือฉบับสมบูรณ์ระบบควบคุมโดยผู้ปกครอง}

ก่อนที่จะมอบอุปกรณ์ให้กับลูก สิ่งแรกที่ต้องทำไม่ใช่การเปิดเครื่อง แต่คือการตั้งค่าระบบควบคุมโดยผู้ปกครอง ทั้งสองแพลตฟอร์มหลักมีระบบในตัวให้เลือกใช้:

\booksubsection{iOS: เวลาใช้งานหน้าจอ}

\begin{itemize}
  \item \textbf{ขีดจำกัด App}:ตั้งค่า > เวลาใช้งานหน้าจอ > ขีดจำกัด App สามารถกำหนดระยะเวลาการใช้งานรายวันสำหรับแต่ละหมวดหมู่ App ได้ (เช่น เกม 30 นาที)
  \item \textbf{เวลาพักการใช้งาน}:กำหนด "ช่วงเวลาดับไฟ" ประจำวัน (เช่น 22:00-7:00) ในช่วงเวลาดังกล่าว App ทั้งหมดจะถูกปิดใช้งานและกลายเป็นสีเทา เหลือเพียงโทรศัพท์ให้ใช้งานได้
  \item \textbf{ข้อจำกัดเนื้อหา}:กรองภาพยนตร์ หน้าเว็บ และการติดตั้ง App ตามช่วงอายุ รวมถึงปิดการซื้อภายใน App
  \item \textbf{Family Sharing}:ใช้ iPhone ของผู้ปกครองจัดการอุปกรณ์ของลูกจากระยะไกล อนุมัติคำขอดาวน์โหลด และดูรายงานการใช้งาน
\end{itemize}

\booksubsection{Android: Google Family Link / ระบบควบคุมโดยผู้ปกครองของแบรนด์}

\begin{itemize}
  \item \textbf{Google Family Link}:โซลูชันอย่างเป็นทางการจาก Google สามารถจัดการระยะเวลาการใช้งานอุปกรณ์ ล็อกอุปกรณ์ และดูรายงานกิจกรรมจากระยะไกลได้ สำหรับ Android ในประเทศจีน ขึ้นอยู่กับความพร้อมใช้งานของแต่ละรุ่น
  \item \textbf{พื้นที่สำหรับเด็ก Xiaomi}:ตั้งค่า > พื้นที่สำหรับเด็ก สามารถกำหนดระยะเวลาการใช้งาน และ App ในรายการที่อนุญาต การออกจากระบบต้องใช้รหัสผ่านของผู้ปกครอง
  \item \textbf{โหมดนักเรียน Huawei}:ตั้งค่า > การใช้งานมือถือเพื่อสุขภาพ > โหมดนักเรียน จำกัดช่วงเวลาการใช้งานและ App รองรับการจัดการระยะไกล
  \item \textbf{โหมดเด็ก OPPO/vivo}:เปิดใช้งานได้ใน "การตั้งค่า" หรือ App "พื้นที่สำหรับเด็ก" เฉพาะเจาะจง หลักการทำงานคล้ายคลึงกัน
\end{itemize}

\begin{center}
\begin{tabularx}{\linewidth}{|X|X|X|}
\hline
\textbf{ฟังก์ชัน} & \textbf{iOS เวลาใช้งานหน้าจอ} & \textbf{Android Family Link / โซลูชันของแบรนด์} \tabularnewline
\hline
ขีดจำกัด App & รองรับ แบ่งตามหมวดหมู่ & รองรับ แบ่งตาม App \tabularnewline
เวลาพักการใช้งาน & รองรับ & รองรับ (แตกต่างกันเล็กน้อยตามแบรนด์) \tabularnewline
การกรองเนื้อหา & ระดับระบบ ค่อนข้างทั่วถึง & ขึ้นอยู่กับแบรนด์ มาตรฐานไม่เท่ากัน \tabularnewline
การจัดการระยะไกล & Family Sharing & Family Link / คลาวด์ของแบรนด์ \tabularnewline
การอนุมัติการติดตั้ง App & ต้องได้รับอนุมัติจากผู้ปกครอง & รองรับโดย Family Link \tabularnewline
\hline
\end{tabularx}
\end{center}

\begin{tcbwarn}{การควบคุมโดยผู้ปกครองไม่ใช่เรื่อง "ทำครั้งเดียวจบ"}
เด็กๆ มักจะหาวิธีบายพาสได้เสมอ (ค้นหาในอินเทอร์เน็ต “วิธีแฮ็กเวลาใช้งานหน้าจอ” จะเจอบทแนะนำเพียบ) มาตรการทางเทคนิคเป็นเพียงขั้นต่ำ แนวป้องกันที่แท้จริงคือการสื่อสารและความไว้วางใจระหว่างผู้ปกครองกับลูก แนะนำให้ทบทวนรายงานการใช้งานร่วมกับลูกทุกเดือน แทนที่จะคอยเฝ้าระวังฝ่ายเดียว

\end{tcbwarn}

\booksection{โหมดสำหรับวัยรุ่น}

แพลตฟอร์มเนื้อหาหลักต่างมี “โหมดสำหรับวัยรุ่น” ในตัว เมื่อเปิดใช้งานจะจำกัดเวลาใช้งาน กรองเนื้อหาที่ไม่เหมาะสม และปิดการเติมเงินหรือให้ทิป นี่เป็นฟังก์ชันบังคับที่กฎหมายของรัฐกำหนด แต่\textbf{ต้องให้ผู้ปกครองเป็นฝ่ายเปิดใช้งานและตั้งรหัสผ่านเอง}:

\begin{itemize}
  \item \textbf{โหมดสำหรับวัยรุ่นของ Douyin/Kuaishou}:จำกัด 40 นาทีต่อวัน (ห้ามใช้งาน 22:00-6:00) เนื้อหาที่แสดงจะคัดเฉพาะวิดีโอประเภทวิทยาศาสตร์และการศึกษาเท่านั้น ไม่สามารถรับชมสด ให้ทิป หรือเติมเงินได้
  \item \textbf{โหมดสำหรับวัยรุ่นของ WeChat}:จำกัด WeChat Pay (ต้องมีการอนุญาตวงเงินจากผู้ปกครอง) จำกัดเนื้อหาในบัญชีสาธารณะและช่องวิดีโอ ฟังก์ชันเช่นเขย่าและคนใกล้เคียงจะถูกปิดเป็นค่าเริ่มต้น
  \item \textbf{ระบบป้องกันการเสพติดเกม}:การยืนยันตัวตนและการจำกัดเวลาระดับชาติ ผู้เยาว์สามารถเข้าสู่ระบบเกมได้เฉพาะวันศุกร์ วันเสาร์ วันอาทิตย์ และวันหยุดนักขัตฤกษ์ เวลา 20:00-21:00 น. เท่านั้น วันละ 1 ชั่วโมง
\end{itemize}

\textbf{อย่าใช้บัตรประชาชนของผู้ปกครองเพื่อ “ยืมบัญชี”}\par

เด็กหลายคนใช้บัตรประชาชนของผู้ปกครองลงทะเบียนบัญชีเกมเพื่อบายพาสระบบป้องกันการเสพติด ซึ่งนอกจากจะละเมิดข้อตกลงผู้ใช้แล้ว ยังทำให้เด็กสูญเสียการปกป้องอีกด้วย ผู้ปกครองควรเก็บรหัสผ่านการชำระเงินและข้อมูลประจำตัวของตนเองไว้ให้ดี อย่า “ยืมบัญชี” ให้ลูกเพียงเพื่อความสะดวก

\booksection{การให้ความรู้ด้านการใช้อินเทอร์เน็ต}

เทคโนโลยีสามารถบล็อกความเสี่ยงได้ 90\% ส่วน 10\% ที่เหลือขึ้นอยู่กับวิจารณญาณของเด็ก แทนที่จะบอกว่า “ห้ามเล่นอินเทอร์เน็ต” สิ่งที่สำคัญกว่าคือการสอนพวกเขาว่า “จะเล่นอินเทอร์เน็ตอย่างไร”:

\begin{itemize}
  \item \textbf{การปกป้องความเป็นส่วนตัว}:ไม่เปิดเผยชื่อจริง โรงเรียน ที่อยู่บ้านให้คนแปลกหน้า ไม่ส่งรูปส่วนตัว บัตรประจำตัว รูปครอบครัว ปิดฟังก์ชั่นระบุตำแหน่งเป็นค่าเริ่มต้น เปิดใช้เฉพาะเมื่อจำเป็นต้องใช้งาน
  \item \textbf{การรับรู้และการรับมือกับการรังแกทางไซเบอร์}:การถูกด่าทอ การกีดกัน การดัดแปลงรูปภาพอย่างมุ่งร้าย ล้วนถือเป็นการรังแก บอกลูกว่าหากพบเจออย่าไปตอบโต้ด้วยคำหยาบคาย อย่าลบหลักฐาน และให้บอกผู้ปกครองทันที ใน WeChat สามารถกดค้างที่ข้อความเพื่อร้องเรียนได้ แพลตฟอร์มจะดำเนินการกับบัญชีผู้กระทำ
  \item \textbf{การวิเคราะห์เนื้อหาอย่างมีวิจารณญาณ}:สอนลูกให้แยกแยะระหว่าง "ข้อเท็จจริง" และ "ความคิดเห็น" จดจำได้ว่าอัลกอริทึมจะแสดงเฉพาะสิ่งที่อยากเห็น ซึ่งไม่ได้หมายความว่าจะเป็นภาพรวมของโลก หากเจอโฆษณาแบบ "ได้ของฟรี" ให้ไม่เชื่อทั้งหมด
  \item \textbf{ความปลอดภัยทางอินเทอร์เน็ต}:ไม่คลิกลิงก์แปลกหน้า ไม่สแกน QR Code ที่ไม่ทราบที่มา ไม่ดาวน์โหลด App นอก Store ที่ไม่เป็นทางการ ไม่แชร์รหัสผ่านให้ใคร (รวมถึง "เพื่อนสนิท") หากมีใครขอรหัสยืนยันให้ถือว่าเป็นมิจฉาชีพทั้งหมด
\end{itemize}

\begin{tcbpractice}{ลงมือทำ: ทำสัญญาฉบับหนึ่งเรื่อง "กติกาการใช้มือถือ" ร่วมกับลูก}
การเขียนกติกาลงไป โดยให้ทั้งสองฝ่ายเซ็นชื่อและติดไว้ในที่ที่เห็นได้ชัดเจน จะได้ผลดีกว่าการตกลงกันปากต่อปาก ขอแนะนำให้ระบุสิ่งต่อไปนี้:

\begin{itemize}
  \item ช่วงเวลาใช้งาน: เช่น "ใช้ได้หลังทำการบ้านเสร็จ 19:30-20:30 น."
  \item พื้นที่ห้ามใช้: บนโต๊ะอาหาร ก่อนนอนในห้องนอน และในห้องเรียน ให้เก็บมือถือทุกครั้ง
  \item ขอบเขตของเนื้อหา: ไม่ติดตั้ง App ที่ไม่ได้แจ้งผู้ปกครอง ไม่คุยส่วนตัวกับคนแปลกหน้า
  \item ผลที่ตามมาเมื่อผิดกติกา: เช่น "เกินเวลา 1 ครั้ง เวลาใช้งานของวันถัดไปจะถูกหักครึ่งหนึ่ง"
  \item สัญญาของผู้ปกครอง: จะไม่ "เล่นมือถือ" ต่อหน้าลูกเกินเวลาที่ตกลงกันไว้ เพื่อเป็นแบบอย่างที่ดี
\end{itemize}

\end{tcbpractice}

\booksection{พฤติกรรมการใช้งานเพื่อสุขภาพ}

ดวงตา กระดูกสันหลัง และการนอนหลับของเด็กยังอยู่ในช่วงพัฒนาการ ผลกระทบจากมือถือต่อร่างกายจึงรุนแรงกว่าผู้ใหญ่:

\begin{itemize}
  \item \textbf{คำแนะนำเรื่องเวลาหน้าจอ}:อายุ 3-6 ปี ไม่เกิน 1 ชั่วโมงต่อวัน อายุ 6-12 ปี ประเภทการเรียนรู้ไม่จำกัด ประเภทบันเทิงไม่เกิน 1 ชั่วโมง อายุ 12 ปีขึ้นไป ให้ตกลงกับลูก แต่ประเภทบันเทิงต้องไม่เกิน 2 ชั่วโมง
  \item \textbf{พฤติกรรมการดูแลสายตา}:ปฏิบัติตาม "กฎ 20-20-20" — มองหน้าจอทุก 20 นาที ให้เงยหน้ามองไกล 20 ฟุต (ประมาณ 6 เมตร) เป็นเวลา 20 วินาที ไม่เล่นมือถือ 1 ชั่วโมงก่อนนอน เพราะแสงสีฟ้าจะยับยั้งเมลาโทนินส่งผลกระทบต่อการนอนหลับ
  \item \textbf{การจัดท่าทาง}:จอมือถือควรอยู่ระดับเดียวกับสายตา ไม่ก้มหน้าหรือไหล่คู่ ใช้งานต่อเนื่องไม่เกิน 30 นาที แล้วควรลุกขึ้นมาขยับร่างกาย
  \item \textbf{กติการ่วมกันในครอบครัว}:กำหนด "พื้นที่ปลอดมือถือ" ที่โต๊ะอาหาร ห้องนอน และห้องทำงาน เมื่อดูหนังร่วมกันทั้งครอบครัวหรือทำกิจกรรมกลางแจ้ง ให้เก็บมือถือทุกคน ทำให้ "การวางมือถือ" เป็นวัฒนธรรมของครอบครัวแทนที่จะเป็นการลงโทษเด็ก
\end{itemize}

\begin{tcbinfo}{ผู้ปกครองคือ "ไฟร์วอลล์" ที่ดีที่สุด}
งานวิจัยพบว่า เวลาที่ผู้ปกครองใช้มือถือมีความสัมพันธ์เชิงบวกกับเวลาที่เด็กใช้มือถืออย่างมาก หากสั่งลูกให้ไม่ดูมือถือ แต่ตัวเองกลับเล่นคลิปสั้นจนดึก กติกาก็จะใช้ไม่ได้ผล หากต้องการเปลี่ยนแปลงเด็ก ให้เริ่มจากการวางมือถือของตัวเองก่อน

\end{tcbinfo}

\booksection{มรดกดิจิทัลและเด็ก}

ส่วนนี้มักถูกมองข้าม แต่สำคัญมาก สินทรัพย์ดิจิทัลของเด็ก (บัญชี, รูปภาพ, ยอดเงินคงเหลือ, ไอเทมเกม) ก็ต้องเผชิญกับปัญหาการสืบทอดและความเป็นส่วนตัวเช่นกัน:

\begin{itemize}
  \item \textbf{การสืบทอดสินทรัพย์ดิจิทัลของเด็ก}: หากเด็กเสียชีวิต กระบวนการสืบทอดบัญชีโซเชียล บัญชีเกม และรูปภาพบนคลาวด์จะเหมือนกับผู้ใหญ่ โดยต้องแสดงใบมรณบัตรและหลักฐานการเป็นญาติ ขอแนะนำให้ผู้ปกครองตั้งค่า "ผู้ติดต่อมรดก" หรือ "ผู้ติดต่อฉุกเฉิน" สำหรับบัญชีสำคัญของเด็กไว้ตั้งแต่เนิ่นๆ (เช่น ผู้ติดต่อมรดกของ Apple, การสืบทอดบัญชี WeChat)
  \item \textbf{ความเป็นส่วนตัวดิจิทัลของผู้เยาว์}: ตามกฎหมายคุ้มครองข้อมูลส่วนบุคคล การประมวลผลข้อมูลส่วนบุคคลของผู้เยาว์ที่มีอายุต่ำกว่า 14 ปี ต้องได้รับความยินยอมจากผู้ปกครอง และต้องมีกฎความเป็นส่วนตัวเฉพาะ ผู้ปกครองควรตรวจสอบนโยบายความเป็นส่วนตัวของ App ที่เด็กใช้บ่อยเป็นประจำ และปิดสวิตช์ "คำแนะนำเฉพาะบุคคล" และ "การแชร์ข้อมูล"
  \item \textbf{ข้อเสนอแนะในการสมัครบัญชี}: บัญชีที่สมัครให้เด็กควรใช้อีเมล/เบอร์โทรศัพท์ที่ผู้ปกครองควบคุมได้ เพื่อให้ง่ายต่อการจัดการรวมศูนย์ในช่วงที่ยังเป็นผู้เยาว์ เมื่อเด็กบรรลุนิติภาวะแล้วจึงค่อยโอนบัญชีให้
\end{itemize}

\begin{tcbwarn}{ระวังอีกด้านของ "มรดกดิจิทัลของเด็ก"}
เด็กๆ ถูก “บันทึก” ตั้งแต่เกิด—รูปเติบโตที่ผู้ปกครองโพสต์, วิดีโอจากโรงเรียนอนุบาล, ข้อมูลสุขภาพจากสมาร์ทวอทช์ เมื่อข้อมูลเหล่านี้ถูกอัปโหลดไปแล้ว จะลบทิ้งให้หมดจดได้ยาก ก่อนโพสต์ ลองคิดดูสักนิด: สิบปีข้างหน้าลูกจะอายที่รูปนี้ไหม? การเคารพตัวตนดิจิทัลของเด็ก เริ่มจากการรู้จักควบคุมตัวเองในการ “โพสต์รูปลูก”

\end{tcbwarn}

\begin{tcbsummary}{สรุปบทนี้}
\begin{itemize}
  \item ไม่มีคำตอบตายตัวว่าอายุเท่าไหร่ถึงควรได้รับมือถือ แต่อายุ 3-6 ปีควรหลีกเลี่ยง, ก่อน 13 ปีให้ใช้สมาร์ทวอทช์หรือโทรศัพท์พื้นฐานไปก่อน
  \item "พอใช้ก็ดีแล้ว": มือถือเครื่องแรกเลือกเครื่องราคาพันกว่าแทนที่จะเป็นเรือธง สเปกเกินความจำเป็นเท่ากับการล่อใจที่เกินความจำเป็น
  \item iOS หน้าจอเวลาใช้งาน และ Android Family Link/การควบคุมโดยผู้ปกครองของแต่ละแบรนด์ เป็นสองระบบป้องกันทางเทคนิค หากใช้ร่วมกันจะมั่นใจได้มากขึ้น
  \item โหมดวัยรุ่นและระบบป้องกันการเสพติดเกมเป็นข้อกำหนดตามกฎหมาย แต่ผู้ปกครองต้องเปิดใช้งานและรักษารหัสผ่านเอง
  \item ทักษะดิจิทัลสำคัญกว่าการบล็อกทางเทคนิค: สอนเด็กให้รู้จักการถูกกลั่นแกล้ง, ปกป้องความเป็นส่วนตัว, และวิเคราะห์เนื้อหาอย่างมีวิจารณญาณ
  \item นิสัยที่ดีต้องเริ่มจากผู้ปกครอง—คุณวางมือถือ ลูกถึงจะวางมือถือตาม
  \item มรดกดิจิทัลและความเป็นส่วนตัวของเด็กเป็นประเด็นใหม่ ควรตั้งค่า "ผู้ติดต่อมรดก" ไว้ตั้งแต่เนิ่นๆ และควบคุมการโพสต์รูปลูก
\end{itemize}

\end{tcbsummary}
